\subsection{Controlled Iteration Design}
The iterations examine three preparation decisions: the scaling method, the feature-selection method and the number of selected features. Each comparison changes one decision while keeping the classifier settings, five holdout seeds (42, 7, 21, 84 and 123), 70/30 stratified split and decision threshold of 0.5 unchanged. Imputation, selection and scaling are fitted within the training pipeline, including each training-only cross-validation fold. The held-out records are used for evaluation only. The completed comparisons below use the retained Python workflow; the native Spark execution status is documented separately so that results from the two implementations are not conflated. A fourth iteration changes the assessment strategy to nested LOPO with three-fold training CV, using the observed completed Research Paper outputs in Section~\ref{subsec:validation-strategy-iteration}.

\subsection{Iteration 1: Change in Scaling Method}
The initial configuration applies a Yeo--Johnson transformation followed by standardization to continuous predictors. Yeo--Johnson is intended to reduce distributional skewness, while standardization puts continuous variables on a common scale. The revised configuration removes the power transformation and uses standardization alone, $z=(x-\mu_{\mathrm{train}})/s_{\mathrm{train}}$. This tests whether the additional distributional transformation contributes useful predictive information or whether the simpler scaling procedure is sufficient. Existing ordinal category codes are passed through without one-hot encoding or continuous-variable scaling.

To isolate the scaling effect, CART selection retains twenty fields in both configurations (Iterations A and B). A second matched comparison repeats the scaling change with RFE retaining twenty fields (Iterations C and D). The same five predictive families are evaluated in each configuration, rather than choosing a scaling method from a single model's result.

\begin{table}[H]
\centering\small
\caption{Controlled fixed-model preparation comparisons: mean test AUC over five matched seeds.}
\label{tab:holdout-preparation}
\begin{tabular}{lrrrr}
\toprule
Model & CART+YJ & CART+Std & RFE+YJ & RFE+Std \\
\midrule
Logistic regression & \textbf{0.861} & 0.836 & \textbf{0.864} & \textbf{0.832} \\
Random forest & 0.810 & 0.811 & 0.808 & 0.808 \\
XGBoost & 0.809 & 0.809 & 0.801 & 0.801 \\
LightGBM & 0.805 & 0.814 & 0.794 & 0.799 \\
Stacking & 0.845 & \textbf{0.839} & 0.837 & 0.824 \\
\bottomrule
\end{tabular}
\end{table}


With CART20, changing from Yeo--Johnson plus standardization to standardization alone changes mean test AUC from 0.861 to 0.836 for logistic regression and from 0.845 to 0.839 for stacking. Random forest and XGBoost change little, while LightGBM increases from 0.805 to 0.814. With RFE20, the corresponding logistic regression and stacking AUCs decrease from 0.864 to 0.832 and from 0.837 to 0.824. These results show that the effect depends on the model; standardization alone does not uniformly improve discrimination. It is retained as the common setting for subsequent selector and feature-count comparisons, and it aligns with the available native Spark baseline. This is a controlled experimental choice, not a claim that it is the best scaler for every family.

\subsection{Iteration 2: Change in Feature-selection Method}
The next change replaces a single CART importance ranking with recursive feature elimination (RFE), while retaining twenty source fields. CART selects the twenty highest-importance predictors from one fitted decision tree. RFE repeatedly fits the decision-tree estimator, removes the least-important remaining predictor and refits until twenty predictors remain. Refitting allows rankings to change after a predictor is removed, but requires more fitting work. Both prototype selectors use the same tree controls (maximum depth five and minimum leaf size four) and operate after training-only imputation. RFE therefore changes the selection procedure, rather than introducing a different ranking estimator.

The principal matched comparison is CART20 plus standardization (Iteration B) versus RFE20 plus standardization (Iteration D). The Yeo--Johnson configurations (A versus C) provide an additional check under the other scaling setting. In the standardization comparison, mean AUC changes from 0.836 to 0.832 for logistic regression, 0.811 to 0.808 for random forest, 0.809 to 0.801 for XGBoost, 0.814 to 0.799 for LightGBM and 0.839 to 0.824 for stacking. RFE20 consequently does not outperform CART20 on mean AUC in this configuration. Under Yeo--Johnson, logistic regression increases slightly from 0.861 to 0.864, while the other families decrease. RFE is retained for the following dimensionality experiment to test the effect of reducing its selected subset; its use is not justified by an unsupported claim of universal superiority.

\subsection{Iteration 3: Change in Number of Selected Features}
The final preparation change reduces the RFE subset from twenty to ten source fields (Iteration D versus Iteration F). Standardization, the decision-tree ranking estimator, fixed classifier parameters and all five holdout memberships remain unchanged. This comparison asks whether a smaller input set preserves predictive performance while reducing the number of required fields. It compares fixed models on both sides, so the feature-count effect is not mixed with the separate Optuna tuning experiment. In stacking, selection is fitted separately inside each base learner's training pipeline; saved artifacts confirm ten retained fields for every base learner in the new run.

\begin{table}[H]
\centering\small
\caption{Completed Python prototype feature-count iteration: fixed settings and matched five-seed holdouts, RFE10 minus RFE20. Both use standardization.}
\label{tab:rfe-feature-count}
\begin{tabular}{lrrrr}
\toprule
Model & RFE20 AUC & RFE10 AUC & $\Delta$ AUC & $\Delta$ sensitivity \\
\midrule
LightGBM & 0.799 $\pm$ 0.027 & 0.804 $\pm$ 0.023 & \textbf{+0.005} & \textbf{+0.000} \\
Logistic regression & \textbf{0.832} $\pm$ 0.042 & 0.807 $\pm$ 0.069 & -0.026 & -0.017 \\
Random forest & 0.808 $\pm$ 0.047 & 0.812 $\pm$ 0.047 & +0.004 & -0.004 \\
Stacking & 0.824 $\pm$ 0.029 & \textbf{0.819} $\pm$ 0.047 & -0.006 & \textbf{+0.000} \\
XGBoost & 0.801 $\pm$ 0.030 & 0.802 $\pm$ 0.038 & +0.001 & \textbf{+0.000} \\
\bottomrule
\end{tabular}
\end{table}
This is a controlled feature-count comparison within the Python implementation. It does not supply native Spark CART/RFE scores. All saved selectors and stacking base pipelines retain exactly ten fields in the new run.


The ten-field configuration produces small mean AUC increases for LightGBM, random forest and XGBoost (approximately 0.005, 0.004 and 0.001). Logistic regression decreases by approximately 0.026 and stacking by 0.006. Sensitivity is unchanged for LightGBM, XGBoost and stacking, but decreases by approximately 0.017 for logistic regression and 0.004 for random forest. Thus, reducing the input count has a model-dependent cost. The small positive AUC differences are descriptive results from five matched holdouts, not evidence of statistical superiority. Both feature counts remain documented rather than selecting ten fields solely because the representation is smaller.

\subsection{Iteration 4: Nested LOPO with Three-fold Training CV}
\label{subsec:validation-strategy-iteration}
This iteration changes the evaluation strategy while preserving the distinction between development and assessment. The reported Iteration 4 results use five stratified 70/30 holdouts with three-fold CV inside the training portion. The earlier 70/15/15 design and LOPO discussion retained from Iteration 3 in Part 4 provide historical context; they are not the split membership used for the current 223/96 holdout result tables. A fixed validation split is easy to interpret, but it allocates fewer records to fitting and makes model choice sensitive to a single small validation set. The repeated-holdout/inner-CV design instead uses the training sample more efficiently while keeping an outer test partition.

\begin{table}[H]
\centering\small
\caption{Validation-strategy iteration: historical design, current completed evaluation and separate research extension.}
\label{tab:validation-strategy-iteration}
\begin{tabularx}{\textwidth}{@{}p{2.6cm}XX@{}}
\toprule
Strategy & Development and assessment roles & Evidence/status \\
\midrule
Fixed 70/15/15 & Training fits the model; a validation partition selects settings; a separate test partition assesses the choice. & Historical Iteration 3 design retained in Part 4; not the current holdout results. \\
Five 70/30 holdouts with inner CV & Each 223-record training set supplies three-fold validation; refit on 223 and test on 96. & Completed current non-Spark and native Spark evidence in Section 7. \\
Nested LOPO with three-fold CV & Exclude one patient; tune on the remaining 318 records with three-fold training CV, then refit and predict the excluded patient. & Completed non-Spark Optuna evidence is reported below; the newer Spark/CUDA workflow remains a separate execution. \\
\bottomrule
\end{tabularx}
\end{table}

\subsubsection{Why Test Leave-One-Patient-Out?}
LOPO is the project name for leave-one-out CV when one record represents one patient. With 319 records, the outer loop has 319 held-out predictions per candidate and seed. Each model can use 318 records for fitting, which tests a different training-size and patient-exclusion scenario from the 223-record holdout training sets. This can reveal sensitivity to individual observations while retaining nearly all observations for training, at a substantial fitting cost. It is not evidence that a model generalises to a new hospital or time period. If repeated observations per patient become available, the split must exclude every record from that patient as a group.

\subsubsection{Observed LOPO Results with Three-fold Training CV}
The Research Paper archive contains completed non-Spark nested Optuna runs for logistic regression, random forest, XGBoost, LightGBM and stacking. Each uses seed 42, twenty Optuna trials per outer patient exclusion, and three-fold validation within the remaining 318 records. An inner split fits on 212 records and validates on 106; Optuna selects settings from the three validation AUCs, then the selected pipeline refits on all 318 records before predicting the excluded patient. Every patient is assessed once, giving 319 held-out predictions per model and 1,595 predictions across the five families. The result table therefore evaluates the complete 319-patient cohort, not one patient: 318 is the training-set size of each outer fit, not the number of test patients or the number of outer folds. The imported complete flags, fold checkpoints and all 31,900 completed trial rows were checked; pooled metrics were independently recomputed from the saved probabilities. The method and evidence are copied into \path{iteration4_BDAS/LOPO} without rerunning training.

\begin{table}[H]
\centering\small
\caption{Observed Research Paper results: nested Optuna LOPO with three-fold training CV, seed 42. Scores pool all 319 held-out patients; each outer model trains on 318 records. Highest displayed scores are bold. No bootstrap estimates are included.}
\label{tab:research-optuna-lopo}
\resizebox{\textwidth}{!}{\begin{tabular}{lrrrrrr}
\toprule
Model & AUC & Accuracy & Precision & Recall & Specificity & F1 \\
\midrule
Logistic regression & \textbf{0.868} & \textbf{0.787} & 0.769 & \textbf{0.826} & 0.747 & \textbf{0.796} \\
Random forest & 0.842 & 0.774 & 0.760 & 0.807 & 0.741 & 0.783 \\
XGBoost & 0.851 & 0.784 & 0.764 & \textbf{0.826} & 0.741 & 0.794 \\
LightGBM & 0.840 & \textbf{0.787} & \textbf{0.775} & 0.814 & \textbf{0.759} & 0.794 \\
Stacking & 0.862 & 0.784 & 0.774 & 0.807 & \textbf{0.759} & 0.790 \\
\bottomrule
\end{tabular}}
\end{table}


The table copies the observed \path{metrics.csv} values from the five completed Research Paper runs. Each value is checked against its saved held-out probabilities and original configuration. The final test metrics pool all 319 excluded-patient predictions; they are not training scores on the 318 records, one-patient metrics or three-fold CV tuning scores. Three-fold CV is used only inside each outer training set to select hyperparameters. No bootstrap resampling, additional fitting or newly estimated scores are included in this evidence.

Only seed 42 has a complete archived run for each family. An across-seed mean or variance cannot be measured from this single run, and neither is fabricated. The table reports the original pooled AUC, accuracy, precision, recall, specificity and F1 exactly as observed, rounded for display. Additional independently executed seeds would be needed to report variation across complete repeated runs.

For these observed results, logistic regression has the highest AUC (0.868), followed by stacking (0.862). Logistic regression and LightGBM tie on accuracy (0.787), LightGBM has the highest precision (0.775), and logistic regression and XGBoost tie on recall (0.826). All five meet the AUC and recall benchmarks of 0.80 in this single complete LOPO evaluation. The different metric leaders reinforce the need to assess discrimination, missed cases and false positives together rather than declaring an unconditional winner.

The copied preprocessing uses training-only median/mode imputation, CART-based RFE selecting twenty \emph{source} fields, Yeo--Johnson followed by standardization for continuous fields, and one-hot encoding of selected categorical fields. These are the historical Research Paper settings; they differ from the ordinal-preserving, standardization-only holdout workflow. The saved stacking search also tunes only the logistic meta-learner's $C$, leaving the base learners fixed, whereas the current holdout stacking search tunes the boosting base learners. These differences are preserved to match the measured output. The LOPO scores therefore do not isolate an effect of changing validation strategy: training size, encoding, transformation and some tuning choices also differ. They do not establish a native Spark or GPU result.

The \path{LOPO} folder retains the original three modelling modules, raw input, requirements, tests and per-model configurations, predictions, trial traces and complete flags. \path{LOPO/results/metrics_combined.csv} copies the observed pooled metrics; \path{LOPO/report_results.py} verifies the saved outputs and regenerates the result table; \path{LOPO/provenance.json} records source/output hashes and validation. The original Research Paper folder remains unchanged. New experiments use a separate output directory so this completed evidence snapshot remains identifiable.

\subsubsection{Training, Tuning and Reporting Rules}
Each outer patient's features and outcome are excluded from fitting, preprocessing and parameter selection. Within the remaining records, stratified CV chooses hyperparameters; selection and scaling are refitted inside each inner training fold. The selected procedure is then refitted on the outer training set and predicts the held-out patient once. The completed non-Spark Optuna evidence above is separate from both the historical fixed-model LOPO archive and the newer nested Spark/CUDA implementation; neither archive establishes execution of the newer Spark/CUDA experiment.

ROC--AUC is calculated from the complete pooled out-of-fold probabilities because a single-patient test fold contains only one outcome class. Accuracy, precision and recall are also calculated from pooled confusion counts for each complete seeded run, rather than averaging undefined one-patient precision/recall values. Multiple seeded runs can report mean and sample standard deviation, but they reuse the same patients and are not independent population replications. Partial checks produce progress and predictions only, not full-evaluation scores.

The Spark research notebook preserves RFE20 or RFE10, cached training-only preparation and resumable prediction/trial evidence. The optional CUDA XGBoost candidate accelerates driver-side boosting while native Spark LR/RF/GBT and RFE remain CPU-based. A fair comparison of strategies should lock the candidate families, preparation, scoring and tuning budget in advance and report the different training sizes and computational costs. Holdout and LOPO outputs remain separately labelled; changing the assessment scheme alone does not demonstrate an improvement.

Together, the three preparation comparisons and validation-strategy iteration document the reason for each change, the variables held constant and the observed trade-off. Training-CV algorithm and parameter decisions remain in Section 6, and detailed execution metrics remain in Section 7. These holdout comparisons inform the next CRISP-DM cycle without using test performance as the tuning objective; external validation and further development are addressed in Future Work.
